\section*{第 \ref{ch:g8:circuit-laws} 章——电路中的电流与电压定律}

\begin{solution}{exo:g8:circuit-laws:1}
\omterm{def:g4:series-circuits:series}{串联}电流定律：一条回路上，处处\omterm{def:g8:intensity-ammeter:intensity}{电流强度}相同。
\omterm{def:g5:series-parallel-circuits:parallel}{节点}定律：流入\omterm{def:g5:series-parallel-circuits:parallel}{节点}的电流之和等于流出的电流之和。\omterm{def:g4:series-circuits:series}{串联}\omterm{def:g8:voltage-voltmeter:voltage}{电压}
定律：\omterm{def:g3:first-electric-circuit:battery}{电池}\omterm{def:g8:voltage-voltmeter:voltage}{电压}等于沿回路各用电器分得的\omterm{def:g8:voltage-voltmeter:voltage}{电压}之和。
\omterm{def:g5:series-parallel-circuits:parallel}{并联}\omterm{def:g8:voltage-voltmeter:voltage}{电压}定律：跨接于相同两\omterm{def:g5:series-parallel-circuits:parallel}{节点}的各\omterm{def:g5:series-parallel-circuits:parallel}{支路}，
\omterm{def:g8:voltage-voltmeter:voltage}{电压}相同。
\end{solution}

\begin{solution}{exo:g8:circuit-laws:2}
$0.8 - 0.55 = \qty{0.25}{A}$.
\end{solution}

\begin{solution}{exo:g8:circuit-laws:3}
$9 - 2.5 - 3.8 = \qty{2.7}{V}$.
\end{solution}

\begin{solution}{exo:g8:circuit-laws:4}
各为 \qty{12}{V}，即\omterm{def:g5:series-parallel-circuits:parallel}{并联}\omterm{def:g8:voltage-voltmeter:voltage}{电压}定律。同一定律
也使家庭电路中每个插座的\omterm{def:g8:voltage-voltmeter:voltage}{电压}都相同，均为 \qty{230}{V}，
即\omterm{def:g5:series-parallel-circuits:parallel}{并联}线路各处\omterm{def:g8:voltage-voltmeter:voltage}{电压}相等。
\end{solution}

\begin{solution}{exo:g8:circuit-laws:5}
以 $I_C$ 为例：已知 $I_{\text{电池}}$ 和 $I_B$ 后，
\omterm{def:g5:series-parallel-circuits:parallel}{节点}定律即可确定 $I_C = \qty{0.2}{A}$——接在此处的任何\omterm{def:g8:intensity-ammeter:ammeter}{电流表}都只能证实这一结果，
（同理，$U_C$ 由\omterm{def:g5:series-parallel-circuits:parallel}{并联}\omterm{def:g8:voltage-voltmeter:voltage}{电压}定律在
已知 $U_B$ 后确定）。
\end{solution}

\begin{solution}{exo:g8:circuit-laws:6}
\qty{2.7}{V} 的\omterm{def:g3:first-electric-circuit:bulb}{灯泡}：电流相等时，\omterm{def:g8:voltage-voltmeter:voltage}{电压}分配越大，
说明它对电流的阻碍作用越强。下一章将\omterm{def:g6:measurement-in-science:measuring}{测量}
这种阻碍作用，即电阻。
\end{solution}

\begin{solution}{exo:g8:circuit-laws:7}
$2.2 + 1.9 + 1.4 = \qty{5.5}{V}$，而\omterm{def:g3:first-electric-circuit:battery}{电池}\omterm{def:g8:voltage-voltmeter:voltage}{电压}为 \qty{6}{V}：
相差半伏，超出此类读数的正常误差范围；
必有一个表盘读错了，或有一段\omterm{def:g8:voltage-voltmeter:voltage}{电压}未被测到。
\end{solution}

\begin{solution}{exo:g8:circuit-laws:8}
$0.4 + 0.3 + 0.3 = \qty{1.0}{A}$ 流出，而流入仅 \qty{0.9}{A}，
分叉处“凭空产生”了 \qty{0.1}{A}——
\omterm{def:g5:series-parallel-circuits:parallel}{节点}定律即可判定其有误。
\end{solution}

\begin{solution}{exo:g8:circuit-laws:9}
回路中各用电器的\omterm{def:g8:voltage-voltmeter:voltage}{电压}之和必须等于\omterm{def:g3:first-electric-circuit:battery}{电池}的 \qty{4.5}{V}：
断开的\omterm{ex:g3:first-electric-circuit:switch}{开关} \qty{4.5}{V} $+$ 不亮的\omterm{def:g3:first-electric-circuit:bulb}{灯泡} \qty{0}{V} $= 4.5$——
平衡。（\omterm{ex:g3:first-electric-circuit:switch}{开关}闭合后，\omterm{def:g8:voltage-voltmeter:voltage}{电压}重新分配：\omterm{ex:g3:first-electric-circuit:switch}{开关}两端\omterm{def:g8:voltage-voltmeter:voltage}{电压}接近零，
\omterm{def:g3:first-electric-circuit:bulb}{灯泡}两端\omterm{def:g8:voltage-voltmeter:voltage}{电压}接近 \qty{4.5}{V}，仍然平衡。）
\end{solution}

\begin{solution}{exo:g8:circuit-laws:10}
变化前：$3 \times 1.5 = \qty{4.5}{V}$。变化后：$4.5 - 1.5 =
\qty{3.0}{V}$——一节\omterm{def:g3:first-electric-circuit:battery}{电池}反接，\omterm{def:g8:voltage-voltmeter:voltage}{电压}相减。\omterm{def:g3:first-electric-circuit:bulb}{灯泡}
由 \qty{4.5}{V} 时的亮度变为 \qty{3}{V} 时的较暗亮度。
\end{solution}

\begin{solution}{exo:g8:circuit-laws:11}
\omterm{def:g5:series-parallel-circuits:parallel}{支路}一跨接\omterm{def:g3:first-electric-circuit:battery}{电池}：$U_E + U_F = \qty{6}{V}$（沿\omterm{def:g5:series-parallel-circuits:parallel}{支路}的\omterm{def:g4:series-circuits:series}{串联}\omterm{def:g8:voltage-voltmeter:voltage}{电压}定律），
故 $U_F = 6 - 2.2 =
\qty{3.8}{V}$。\omterm{def:g5:series-parallel-circuits:parallel}{支路}二跨接相同两\omterm{def:g5:series-parallel-circuits:parallel}{节点}：$U_G =
\qty{6}{V}$（\omterm{def:g5:series-parallel-circuits:parallel}{并联}\omterm{def:g8:voltage-voltmeter:voltage}{电压}定律）。
\end{solution}

\begin{solution}{exo:g8:circuit-laws:12}
$I_R = 0.6 - 0.45 = \qty{0.15}{A}$（\omterm{def:g5:series-parallel-circuits:parallel}{节点}定律）。$U_R = U_Q
= \qty{2.3}{V}$（\omterm{def:g5:series-parallel-circuits:parallel}{并联}\omterm{def:g8:voltage-voltmeter:voltage}{电压}定律）。$U_P = 6 - 2.3 =
\qty{3.7}{V}$（\omterm{def:g4:series-circuits:series}{串联}\omterm{def:g8:voltage-voltmeter:voltage}{电压}定律）。一个可被判错的说法：
“$I_R = \qty{0.2}{A}$”——\omterm{def:g5:series-parallel-circuits:parallel}{节点}处的电流对不上：
$0.45 + 0.2 = 0.65 \neq 0.6$；或者“$U_R = \qty{2.6}{V}$”——
\omterm{def:g5:series-parallel-circuits:parallel}{并联}\omterm{def:g5:series-parallel-circuits:parallel}{支路}的\omterm{def:g8:voltage-voltmeter:voltage}{电压}不能互相矛盾。
\end{solution}

\begin{solution}{pb:g8:circuit-laws:1}
\textbf{1.} \omterm{def:g4:series-circuits:series}{串联}。\omterm{def:g4:series-circuits:series}{串联}时两只\omterm{def:g3:first-electric-circuit:bulb}{灯泡}依次阻碍电流，
回路中的电流远小于单只\omterm{def:g3:first-electric-circuit:bulb}{灯泡}的
\qty{0.4}{A}——测得的 \qty{0.2}{A} 与此相符。
（\omterm{def:g5:series-parallel-circuits:parallel}{并联}则是两条完整\omterm{def:g5:series-parallel-circuits:parallel}{支路}的电流\emph{相加}。）
\textbf{2.} \omterm{def:g5:series-parallel-circuits:parallel}{并联}：两条\omterm{def:g5:series-parallel-circuits:parallel}{支路}各约 \qty{0.4}{A}，
\omterm{def:g3:first-electric-circuit:battery}{电池}输出约 \qty{0.8}{A}。
\textbf{3.} \omterm{def:g4:series-circuits:series}{串联}\omterm{def:g8:voltage-voltmeter:voltage}{电压}定律将 \qty{6}{V} 分到
两只相同\omterm{def:g3:first-electric-circuit:bulb}{灯泡}上，每只约 \qty{3}{V}。
\textbf{4.} \omterm{def:g5:series-parallel-circuits:parallel}{并联}——每只\omterm{def:g3:first-electric-circuit:bulb}{灯泡}都将跨接整个
\qty{6}{V}。从观察窗看，\omterm{def:g4:series-circuits:series}{串联}的一对\omterm{def:g3:first-electric-circuit:bulb}{灯泡}发光较暗：
每只\omterm{def:g3:first-electric-circuit:bulb}{灯泡}两端的\omterm{def:g8:voltage-voltmeter:voltage}{电压}只有平常的一半。
\textbf{5.} $I_Z = 0.5 - 0.35 = \qty{0.15}{A}$（\omterm{def:g5:series-parallel-circuits:parallel}{节点}定律）；
$I_X = \qty{0.5}{A}$——X 位于未分叉的通路上（\omterm{def:g4:series-circuits:series}{串联}
电流定律）。
\textbf{6.} $U_{\text{一对}} = 6 - 2.8 = \qty{3.2}{V}$（回路\omterm{def:g8:voltage-voltmeter:voltage}{电压}
定律）；$U_Y = U_Z = \qty{3.2}{V}$（\omterm{def:g5:series-parallel-circuits:parallel}{并联}\omterm{def:g8:voltage-voltmeter:voltage}{电压}定律）。
\textbf{7.} \omterm{def:g8:voltage-voltmeter:voltage}{电压}相同，电流路径不同：两只\omterm{def:g3:first-electric-circuit:bulb}{灯泡}对电流的
阻碍并不相同——在相同\omterm{def:g8:voltage-voltmeter:voltage}{电压}下通过电流较小的 Z，
阻碍作用更强（路更难走）。
\textbf{8.} 变化：去掉 Y 所在\omterm{def:g5:series-parallel-circuits:parallel}{支路}即少了一条通路，
剩下的 Z 通路对电流的阻碍比原来两者共同的阻碍还要大，因此
\omterm{def:g3:first-electric-circuit:battery}{电池}输出的总电流变小。
\textbf{9.} 第一条：回路定律，$2.5 + 4.1 =
\qty{6.6}{V}$，而电源\omterm{def:g8:voltage-voltmeter:voltage}{电压}只有 \qty{6}{V}。第二条：
\omterm{def:g5:series-parallel-circuits:parallel}{节点}定律：流入 \qty{0.3}{A}，流出 $0.2 + 0.15 =
\qty{0.35}{A}$。
\textbf{10.} $6 - 2.5 = \qty{3.5}{V}$。
\textbf{11.} $0.2 + 0.15 = \qty{0.35}{A}$。
\textbf{12.} 因为四条定律对任何真实电路都相互约束：
每个读数都须与其他读数的和相符，因此一个
编造的数字至少会违反一条定律，
仅凭算术即可判定其有误。
\textbf{13.} 例如：“同一条路，同一股电流——\omterm{def:g4:series-circuits:series}{串联}电流定律。
在每个分叉处，流进与流出平衡——\omterm{def:g5:series-parallel-circuits:parallel}{节点}定律。
沿任一回路，\omterm{def:g3:first-electric-circuit:battery}{电池}的“水流”
被各用电器耗尽，一滴不剩——\omterm{def:g4:series-circuits:series}{串联}\omterm{def:g8:voltage-voltmeter:voltage}{电压}定律。
挂在同样两个\omterm{def:g5:series-parallel-circuits:parallel}{节点}之间的各\omterm{def:g5:series-parallel-circuits:parallel}{支路}，
都从同一高度取水——\omterm{def:g5:series-parallel-circuits:parallel}{并联}\omterm{def:g8:voltage-voltmeter:voltage}{电压}定律。”
\end{solution}
